\begin{landscape}
\begin{table}[!ht]
\centering
\caption{Ближайшие классы постановок и область настоящей модели. Значения «да», «частично» и «нет» показывают наличие свойства в классе постановок, а не эмпирическую эквивалентность или новизну.}
\label{tab:closest-prior-work}
\setlength{\tabcolsep}{2.1pt}
\renewcommand{\arraystretch}{1.18}
\renewcommand{\textsuperscript}[1]{\ensuremath{^{\mathrm{#1}}}}
\scriptsize
\begin{tabular}{>{\raggedright\arraybackslash}p{2.85cm}>{\centering\arraybackslash}p{1.45cm}>{\centering\arraybackslash}p{1.9cm}>{\centering\arraybackslash}p{1.65cm}>{\centering\arraybackslash}p{1.7cm}>{\centering\arraybackslash}p{1.6cm}>{\centering\arraybackslash}p{1.45cm}>{\centering\arraybackslash}p{1.55cm}>{\centering\arraybackslash}p{1.95cm}>{\raggedright\arraybackslash}p{2.6cm}>{\raggedright\arraybackslash}p{4.15cm}}
\toprule
Класс постановок & DAG & \shortstack{Потоки\\исполнителей} & \shortstack{Сервис\\мощности 1} & \shortstack{Повторный\\сервис} & \shortstack{Удержание\\потока} & \shortstack{$k$ задачи\\и режима} & \shortstack{Разделение\\$a_i/h_i$} & \shortstack{Издержки\\координации\\и интеграции} & \shortstack{Основной\\результат} & \shortstack{Калибровка реального\\программного процесса} \\
\midrule
RCPSP и многорежимная RCPSP \cite{brucker1999resource}\textsuperscript{a}
& да & частично & да & частично & нет & нет & нет & частично & makespan, стоимость, допустимость & нет \\

Параллельные машины с общим сервером \cite{hall2000parallel,kravchenko1997parallel,elidrissi2023loading,chakhlevitch2009reentrant}\textsuperscript{b}
& нет & да & да & частично & частично & нет & частично & нет & makespan, алгоритмические оценки & нет \\

HRI: число роботов на оператора и распределение внимания \cite{olsen2004fanout,crandall2005validating,crandall2011computing,cummings2008predicting}\textsuperscript{c}
& нет & да & да & да & частично & нет & частично & частично & пропускная способность, выполнение миссии & нет \\

Универсальные MAS и централизованная координация \cite{kim2025scaling}\textsuperscript{d}
& нет & частично & нет & частично & нет & нет & нет & частично & успешность, качество; иногда стоимость & нет \\

Агенты программирования: MAGIS и CAID \cite{tao2024magis,geng2026caid}\textsuperscript{e}
& частично & частично & нет & частично & нет & нет & нет & частично & успешность, качество; CAID также время и стоимость & нет \\

Настоящая работа\textsuperscript{f}
& да & да & да & да & да & да & да & да & makespan принятого результата & нет \\
\bottomrule
\end{tabular}

\vspace{0.5em}
\begin{minipage}{0.97\linewidth}
\scriptsize
\textsuperscript{a} Возобновляемый ресурс RCPSP не обязан представлять одинаковые агентные потоки; повторное обслуживание и издержки требуют явных операций, а стандартная постановка не вводит автоматическое удержание потока в ожидании другого ресурса.

\textsuperscript{b} Повторный сервис присутствует в вариантах с загрузкой/выгрузкой и возвратом, но не в постановках только с подготовкой. Немедленная выгрузка может удерживать машину, тогда как возвратная модель с удалённым сервером освобождает её; машинная и сервисная длительности не откалиброваны как доли работы агента и человека.

\textsuperscript{c} Автономная система ожидает оператора или теряет результативность, однако это не равнозначно удержанию программного потока. Интервалы взаимодействия и автономности не являются программными $a_i/h_i$; очередь, потеря ситуационной осведомлённости и переориентация не включают интеграцию кода.

\textsuperscript{d} Последовательная взаимозависимость измеряется, но явный DAG задач не строится; число агентов или ролей не доказывает перекрытие интервалов. Централизованный координатор не равнозначен наблюдаемому сервису мощности 1, повторные сообщения --- человеческому циклу, а различия результатов --- калибровке издержек M4.

\textsuperscript{e} CAID строит план зависимостей и показывает перекрывающиеся ветви; для MAGIS и класса в целом эти свойства не установлены. Автоматизированные планировщик, контроль качества и доработка не являются человеческим сервисом; ожидание зависимостей или слияния не тождественно локальной блокировке. CAID наблюдает слияние, тестирование, время и стоимость, но не калибрует человеко-агентный процесс. Оценивание на программных наборах задач не является эмпирической калибровкой реального рабочего процесса с наблюдаемыми человеческими фазами.

\textsuperscript{f} В статье используются синтетические сценарии; эмпирическая калибровка и перспективная проверка остаются будущей работой.

\medskip
Статья не заявляет новизной DAG, makespan, общий ресурс, повторный сервис, возврат к исполнителю, удержание ресурса, дуги ресурсного порядка или число автономных исполнителей на одного оператора.
\end{minipage}
\end{table}
\end{landscape}
